\section{Aksesibilitas Tabel dan Form}

Tabel dan form adalah dua elemen HTML yang paling sering menimbulkan masalah aksesibilitas jika tidak diimplementasikan dengan benar. Pengguna pembaca layar menavigasi tabel sel per sel dan mengandalkan header untuk memahami konteks setiap nilai data. Tanpa pengaitan header yang tepat, sebuah angka ``85'' di dalam tabel tidak memberikan informasi apakah itu nilai ujian, usia, atau harga. Demikian pula, form tanpa label yang tepat menjadikan proses pengisian sangat sulit bagi pengguna yang tidak dapat melihat tata letak visual \cite{mdnAccessibility}.

\subsection{Tabel Aksesibel}

Tabel yang aksesibel memerlukan beberapa komponen kunci. Pertama, setiap tabel harus memiliki \code{<caption>} yang mendeskripsikan isi tabel secara ringkas. Kedua, gunakan \code{<th>} untuk semua sel header, bukan \code{<td>} yang dicetak tebal dengan CSS. Ketiga, tambahkan atribut \code{scope="col"} atau \code{scope="row"} pada setiap \code{<th>} untuk menjelaskan arah cakupan header. Keempat, untuk tabel dengan header kompleks atau multi-level, gunakan pasangan \code{id} pada \code{<th>} dan \code{headers} pada \code{<td>}. Kelima, hindari penggunaan \code{colspan} dan \code{rowspan} yang berlebihan karena mempersulit navigasi sel \cite{webdevAccessibility}.

\subsection{Form Aksesibel}

Form yang aksesibel dibangun di atas tiga pilar utama: pelabelan yang tepat, pengelompokan yang logis, dan umpan balik error yang jelas. Setiap kontrol input \textbf{wajib} memiliki label yang terasosiasi---baik melalui \code{<label for="id">} secara eksplisit maupun melalui \code{aria-label} atau \code{aria-labelledby}. Grup field yang terkait (misalnya radio button untuk jenis kelamin) harus dibungkus dengan \code{<fieldset>} dan \code{<legend>} agar pembaca layar menjelaskan konteks pilihan. Pesan error harus dihubungkan dengan input yang bermasalah menggunakan \code{aria-describedby} \cite{mdnAccessibility}.

\begin{table}[htbp]
\centering
\begin{tabular}{|P{3.5cm}|P{3.5cm}|P{5cm}|}
\hline
\textbf{Teknik} & \textbf{Penerapan} & \textbf{Manfaat Aksesibilitas} \\
\hline
\code{<label for>} & Mengaitkan label ke input & Screen reader membaca label saat input difokuskan \\
\hline
\code{<fieldset>} + \code{<legend>} & Mengelompokkan field terkait & Memberikan konteks grup pada radio/checkbox \\
\hline
\code{aria-required} & Menandai field wajib & Screen reader mengumumkan ``required'' \\
\hline
\code{aria-invalid} & Menandai input yang salah & Screen reader mengumumkan ``invalid'' \\
\hline
\code{aria-describedby} & Hubungkan pesan error & Screen reader membaca pesan error \\
\hline
\code{autocomplete} & Petunjuk auto-fill & Mempercepat pengisian form \\
\hline
\end{tabular}
\caption{Teknik aksesibilitas form HTML}
\end{table}

\begin{contoh}
\textbf{Studi Kasus: Form Pendaftaran Mahasiswa yang Aksesibel}

Berikut adalah contoh form pendaftaran yang menerapkan semua prinsip aksesibilitas: setiap input memiliki label, field dikelompokkan dengan fieldset/legend, radio button memiliki konteks, pesan error ditautkan dengan \code{aria-describedby}, dan field wajib ditandai secara visual maupun untuk pembaca layar.
\end{contoh}

\begin{webcode}[language=HTML, caption={Studi kasus: form pendaftaran aksesibel lengkap}]
<form action="/daftar" method="post" novalidate>
  <fieldset>
    <legend>Data Diri</legend>

    <label for="nama">Nama Lengkap
      <span aria-hidden="true">*</span></label>
    <input type="text" id="nama" name="nama"
           required aria-required="true"
           aria-describedby="err-nama"
           autocomplete="name" />
    <span id="err-nama" role="alert"></span>

    <label for="email">Alamat Email
      <span aria-hidden="true">*</span></label>
    <input type="email" id="email" name="email"
           required aria-required="true"
           aria-describedby="err-email"
           autocomplete="email" />
    <span id="err-email" role="alert"></span>

    <label for="pass">Password
      <span aria-hidden="true">*</span></label>
    <input type="password" id="pass" name="pass"
           required minlength="8"
           aria-required="true"
           aria-describedby="hint-pass"
           autocomplete="new-password" />
    <span id="hint-pass">Minimal 8 karakter</span>
  </fieldset>

  <fieldset>
    <legend>Jenis Kelamin</legend>
    <label>
      <input type="radio" name="gender"
             value="L" required /> Laki-laki
    </label>
    <label>
      <input type="radio" name="gender"
             value="P" /> Perempuan
    </label>
  </fieldset>

  <fieldset>
    <legend>Jurusan</legend>
    <label for="jurusan">Pilih Jurusan:</label>
    <select id="jurusan" name="jurusan" required
            aria-required="true">
      <option value="">-- Pilih --</option>
      <option value="if">Teknik Informatika</option>
      <option value="si">Sistem Informasi</option>
    </select>
  </fieldset>

  <label>
    <input type="checkbox" name="agree" required
           aria-required="true" />
    Saya menyetujui syarat dan ketentuan
  </label>

  <button type="submit">Daftar Sekarang</button>
</form>
\end{webcode}

\begin{catatan}
\textbf{Pengujian Aksesibilitas Form:} Untuk memverifikasi aksesibilitas form, lakukan langkah berikut:
\begin{enumerate}
  \item \textbf{Uji keyboard}: pastikan semua field dapat dijangkau dan diisi hanya dengan keyboard (Tab, Shift+Tab, Enter, Space).
  \item \textbf{Uji pembaca layar}: gunakan NVDA (gratis untuk Windows), VoiceOver (macOS/iOS), atau TalkBack (Android) untuk menelusuri form.
  \item \textbf{Periksa fokus visible}: setiap elemen yang difokuskan harus memiliki indikator visual yang jelas (outline).
  \item \textbf{Gunakan alat otomatis}: aXe DevTools, WAVE, atau Lighthouse Accessibility audit untuk mendeteksi masalah umum \cite{webdevAccessibility}.
\end{enumerate}
\end{catatan}

